\subsection{赋值除环}

定义 2. 除环 K 上的绝对值是一个映射 \(x \to |x|\) ，它把 K 映入 \(\mathbf{R}_{+}\) ，且满足下列条件：

\((\mathbf{VM}_{\mathbf{I}})\) \(|x| = 0\) 当且仅当 \(x = 0\) 。

\((\mathbf{VM}_{\mathbf{II}})\quad |xy| = |x|.|y|\) 对一切 \(x,y\) （\(\mathbf{K}\) 中的）成立。

\((\mathbf{VM}_{\mathbf{III}})\quad |x + y|\leqslant |x| + |y|\) 对一切 \(x,y\) （\(\mathbf{K}\) 中的）成立。

由 \((\mathbf{VM}_{\mathbf{II}})\) 有 \(|x| = |\mathrm{I}|.|x|\) ，又由 \((\mathbf{VM}_{\mathbf{I}})\) 至少有一个 \(x\) 使 \(|x| \neq 0\) ，故 \(|\mathrm{I}| = \mathrm{I}\) ；由此推出 \(\mathrm{I} = |- \mathrm{I}|^2\) ，因此 \(|- \mathrm{I}| = \mathrm{I}\) ，从而

\[
| - x | = | - \mathrm{I} | | x | = | x |;
\]

因此 \(|x - y| \leqslant |x| + |y|\) 对一切 \(x, y\) （\(\mathbf{K}\) 中的）成立。于是我们可以说，\(d(x, y) = |x - y|\) 是加法群 \(\mathbf{K}\) 上的不变度量，并且映射 \(x \to |x|\) 是由 K 的非零元素组成的乘法群 \(K^{*}\) 到乘法群 \(R_{+}^{*}\) （由实数 \textgreater0 组成）中的同态。

不变度量 \(|x - y|\) 在 K 上定义了一个度量空间拓扑，它与 K 的加法群结构相容（第 1 号，命题 3）；而且，这一拓扑与 K 的除环结构相容。因为 xy 在 \(K \times K\) 上的连续性由关系式

\[
x y - x _ {0} y _ {0} = (x - x _ {0}) (y - y _ {0}) + (x - x _ {0}) y _ {0} + x _ {0} (y - y _ {0}),
\]

由此得到

\[
\left| x y - x _ {0} y _ {0} \right| \leqslant \left| x - x _ {0} \right|. \left| y - y _ {0} \right| + \left| x _ {0} \right|. \left| y - y _ {0} \right| + \left| y _ {0} \right|. \left| x - x _ {0} \right|.
\]

同样，\(x^{-1}\) 在每一点 \(x_{0} \neq 0\) 处的连续性由恒等式 \(x^{-1} - x_{0}^{-1} = x^{-1} (x_{0} - x)x_{0}^{-1}\) 得出，由 \((\mathrm{VM}_{\mathrm{II}})\) 它给出

\[
\left| x ^ {- 1} - x _ {0} ^ {- 1} \right| = \frac {\left| x - x _ {0} \right|}{\left| x _ {0} \right| . \left| x \right|};
\]

现若 \(\varepsilon > 0\) 满足 \(\varepsilon < |x_0|\) ，则关系 \(|x - x_0| \leqslant \varepsilon\) 蕴含 \(|x| \geqslant |x_0| - \varepsilon\) ，由此得 \(|x^{-1} - x_{0}^{-1}| \leqslant \varepsilon / |x_0| (|x_0| - \varepsilon)\) ；这就建立了 \(x^{-1}\) 在点 \(x_0\) 处的连续性。

定义 3. 赋值除环就是赋有由 K 上一个给定绝对值所定义之结构的除环 K。

赋值除环总被视为赋有由其绝对值定义的拓扑，这使它成为拓扑除环。若 \(K_{0}\) 是赋值除环 K 的一个子除环，则 K 上的绝对值在 \(K_{0}\) 上的限制是 \(K_{0}\) 上的一个绝对值，它在 \(K_{0}\) 上定义由 K 的拓扑诱导的拓扑。

例。1) 设 K 是任意除环。对每个 \(x \in \mathbf{K}\) ，令 \(|x| = \mathbf{i}\) （若 \(x \neq 0\) ），并令 \(|0| = 0\) 。如此定义的映射 \(x \to |x|\) 是 K 上的绝对值，称为平凡绝对值。除环 K 上由绝对值 \(|x|\) 定义的拓扑是离散的，当且仅当 \(|x|\) 是平凡绝对值。此条件显然是充分的；反之，若 K 的拓扑是离散的，则 \(|x|\) 不能取任何 \(\alpha > 0\) 的值，除非它是 \(\mathbf{i}\) ；因为若有 \(|x_0| = \alpha < \mathbf{i}\) ，则序列 \((x_0^n)\) 将由非零项组成且收敛于 \(\mathbf{o}\) ；要处理 \(\alpha > \mathbf{i}\) 的情形，就用 \(x_0^{-1}\) 代替 \(x_0\) 。

\begin{enumerate}
\def\labelenumi{\arabic{enumi})}
\setcounter{enumi}{1}
\item
  实数的绝对值（第 IV 章，§ 1，第 6 号）满足公理 \((\mathrm{VM}_{\mathrm{I}})\) ， \((\mathrm{VM}_{\mathrm{II}})\) 与 \((\mathrm{VM}_{\mathrm{III}})\) ，并且它在域 \(\mathbf{R}\) 上定义的拓扑就是实直线的拓扑。在复数域 \(\mathbf{C}\) （与 \(\mathbf{R}^2\) 等同）上{[}相应地，四元数体 \(\mathbf{H}\) （与 \(\mathbf{R}^4\) 等同）上{]}，欧氏范数同样是绝对值，并且定义域 C（相应地，除环 H）的拓扑（第 VIII 章，§ 1，第 2 号与第 4 号）。
\item
  在除环 K 上，实赋值是定义在 \(K^{*}\) 上、取值于 R 中的函数 v，它满足下列条件：a) 若 \(x \in K^{*}\) 且 \(y \in K^{*}\) ，则 \(v(xy) = v(x) + v(y)\) ；b) 若此外还有 \(x + y \neq 0\) ，则 \(v(x + y) \geqslant \inf(v(x), v(y))\) 。若 a 是任一实数 \textgreater1，则令 \(|x| = a^{-v(x)}\) （对 \(x \neq 0\) ），并令 \(|0| = 0\) ，便可在 K 上定义一个绝对值。因为关系式 \(v(xy) = v(x) + v(y)\) （对 \(x \neq 0\) 与 \(y \neq 0\) ）蕴含对这些 x、y 值的关系式 \(|xy| = |x|.|y|\) ，且当 x、y 之一为零时此关系平凡成立；同样，由关系式 \(v(x + y) \geqslant \inf(v(x), v(y))\) （当 \(x \neq 0\) 、 \(y \neq 0\) 且 \(x + y \neq 0\) 时）推出 \(|x + y| \leqslant \sup(|x|, |y|) \leqslant |x| + |y|\) ，并且当 x、y、 \(x + y\) 之一为零时这些不等式仍然成立。特别地，若 \(v_{p}(x)\) 是有理数域 Q 上的 p 进赋值（x 的素因子分解中 p 的指数），则相应的绝对值 \(|x|_{p} = p^{-v_{p}(x)}\) 称为域 Q 上的 p 进绝对值（参见第 III 章，§ 6，习题 23）。
\end{enumerate}

注。若 \(x\) 是赋值除环中的单位根，则 \(|x| = 1\) ，因为 \(x^n = 1\) 蕴含 \(|x|^n = 1\) ，从而 \(|x| = 1\) 。特别地，有限域上唯一的绝对值是平凡绝对值，因为这种域的每个元素 \(\neq 0\) 者都是单位根。

定义 4. 除环 K 上的两个绝对值称为等价的，如果它们在 K 上定义相同的拓扑。

命题 5. 两个绝对值 \(|x|_1, |x|_2\) （在除环 \(K\) 上，均非平凡绝对值）等价，当且仅当关系 \(|x|_1 < i\) 蕴含 \(|x|_2 < i\) 。这时存在实数 \(\rho > 0\) ，使得 \(|x|_2 = |x|_1^p\) 对一切 \(x \in K\) 成立。

此条件是必要的，因为全体 \(x \in \mathbf{K}\) 中满足 \(|x|_1 < 1\) 者所成的集，与全体 \(x\) 中关于由绝对值 \(|x|_1\) 定义的拓扑使 \(\lim x^n = 0\) 者所成的集相同。

反之，设 \(|x|_1 < \mathrm{i} \Longrightarrow |x|_2 < \mathrm{i}\) 。则有 \(|x|_1 > \mathrm{i} \Longrightarrow |x|_2 > \mathrm{i}\) ，因为 \(|x^{-1}|_1 < \mathrm{i}\) ，从而 \(|x^{-1}|_2 < \mathrm{i}\) 。由假设绝对值 \(|x|_1\) 不是平凡的，故存在 \(x_0 \in \mathbf{K}\) 使 \(|x_0|_1 > \mathrm{i}\) 。令 \(a = |x_0|_1\) ， \(b = |x_0|_2\) ，并令 \(\rho = \log b / \log a > 0\) 。设 \(x \in \mathbf{K}^*\) ，令 \(|x|_1 = |x_0|^{\gamma}\) 。若整数 \(m\) 与 \(n\) 满足 \(n > 0\) 且 \(m/n > \gamma\) ，则 \(|x|_1 < |x_0|_1^{m/n}\) ，因此 \(|x^n x_0^{-m}|_1 < \mathrm{i}\) ；于是 \(|x^n x_0^{-m}|_2 < \mathrm{i}\) ，即 \(|x|_2 < |x_0|_2^{m/n}\) 。类似地，若 \(m/n < \gamma\) ，则得 \(|x|_2 > |x_0|_2^{m/n}\) ；由此推出 \(|x|_2 = |x_0|_2^\gamma\) ；换言之

\[
\log | x | _ {2} = \gamma \log b = \gamma \rho \log a = \rho \log | x | _ {1},
\]

即 \(|x|_2 = |x|_1^p\) 。现在显然，在 \(\mathbf{K}\) 上由 \(|x|_1\) 与 \(|x|_2\) 定义的诸拓扑的零邻域是相同的。

反之，若 \(|x|\) 是 \(K\) 上任一绝对值，则函数 \(|x|^{\circ}\) 是 \(K\) 上的绝对值（与 \(|x|\) 等价），对一切 \(\rho\) （满足 \(0 < \rho \leqslant 1\) ）。只需验证不等式 \(|x + y|^{\circ} \leqslant |x|^{\circ} + |y|^{\circ}\) ；由于 \(|x + y|^{\circ} \leqslant (|x| + |y|)^{\circ}\) ，只需证明：若 \(a > 0\) 且 \(b > 0\) ，则 \((a + b)^{\circ} \leqslant a^{\circ} + b^{\circ}\) 对任何 \(\rho\) （满足 \(0 < \rho \leqslant 1\) ）成立。若令 \(c = a / (a + b)\) 与 \(d = b / (a + b)\) ，则有 \(c + d = 1\) ，待证的不等式为 \(c^{\circ} + d^{\circ} \geqslant 1\) ；但这由关系式 \(c^{\circ} \geqslant c\) 与 \(d^{\circ} \geqslant d\) 立即得出，它们当 \(0 < c \leqslant 1\) 、 \(0 < d \leqslant 1\) 、 \(0 < \rho \leqslant 1\) 时成立。

因此，r \textgreater{} 0 中使 \(|x|^{r}\) 成为绝对值的那些值的集合是 R 中以 o 为左端点的有限或无限区间；若它是有限的，则它显然是闭的；因为若 \(|x + y|^{r} \leqslant |x|^{r} + |y|^{r}\) 对 K 中任意 x、y 及一切满足 \(0 < r < r_{0}\) 的 r 成立，则由连续性该不等式对 \(r = r_{0}\) 仍成立。若 \(|x|^{r}\) 对一切 r \textgreater{} 0 都是绝对值，则我们有

\[
| x + y | \leqslant (| x | ^ {r} + | y | ^ {r}) ^ {1 / r}
\]

对一切 \(x\) 与 \(y\) （\(\mathbf{K}\) 中的）以及一切 \(r > 0\) 成立。现在，若 \(a, b\) 是两个实数 \(\geqslant 0\) ，则有 \(\lim_{r \to \infty} (a^r + b^r)^{1/c} = \sup (a, b)\) ；因为，例如设 \(a \geqslant b\) ，我们有 \(a \leqslant (a^r + b^r)^{1/r} \leqslant 2^{1/r} a\) ，令 \(r \to +\infty\) 即得结果。

于是，若 \(|x|^r\) 对一切 \(r > 0\) 是绝对值，则我们有

\[
| x + y | \leqslant \sup (| x |, | y |)
\]

这可以表述为：\(v(x) = -\log |x| (x \neq 0)\) 是 K 上的赋值。

注。命题 5 的证明表明：若由 \(|x|_2\) 定义的拓扑比由 \(|x|_1\) 定义的拓扑粗，且 \(|x|_1\) 不是平凡的，则 \(|x|_1\) 与 \(|x|_2\) 等价，因为这时关系 \(|x|_2 < 1\) 蕴含 \(|x|_2 < 1\) 。于是，\(K\) 上两个均非平凡的绝对值所定义的拓扑不能比较而不相同。

命题 6. 设 \(\hat{\mathbf{K}}\) 是除环 \(\mathbf{K}\) 赋有绝对值 \(|x|\) 后的完备化，则它是一个除环，且函数 \(|x|\) 可由连续性扩张为 \(\hat{\mathbf{K}}\) 上的绝对值，它定义 \(\hat{\mathbf{K}}\) 的拓扑。

设 \(\mathfrak{F}\) 是 K 上的柯西滤子（关于加法一致结构），它不以 o 为聚点；要证明 \(\hat{\mathbf{K}}\) 是除环，只需证明 \(\mathfrak{F}\) 在映射 \(x \to x^{-1}\) 下的像是柯西滤子基（第 III 章，§ 6，第 8 号，命题 7）。由假设，存在实数 \(\alpha > 0\) 与集 \(\mathrm{A} \in \mathfrak{F}\) ，使得 \(|x| \geqslant \alpha\) 对一切 \(x \in \mathrm{A}\) 成立；另一方面，对每个 \(\varepsilon > 0\) 存在集 \(\mathrm{B} \in \mathfrak{F}\) ，使得 \(\mathrm{B} \subset \mathrm{A}\) 且 \(|x - y| \leqslant \varepsilon\) 对一切 \(x \in \mathrm{B}\)

与 \(y\in \mathbf{B}\) 成立；由此

\[
\left| x ^ {- 1} - y ^ {- 1} \right| = \frac {\left| x - y \right|}{\left| x \right| . \left| y \right|} \leqslant \frac {\varepsilon}{\alpha^ {2}},
\]

这就得出命题的第一部分。不变度量 \(|x - y| = d(x, y)\) 由连续性扩张为 \(\hat{\mathbf{K}}\) 上的度量（§ 2，第 1 号，命题 1），它定义 \(\hat{\mathbf{K}}\) 的拓扑，且由恒等式延拓原理是不变的；我们仍以 \(d(x, y)\) 记这个不变度量。若令 \(|x| = d(0, x)\) （对 \(x \in \hat{\mathbf{K}}\) ），则显然 \(|x|\) 是函数 \(|x|\) （\(\mathbf{K}\) 上的）的连续延拓，从而由恒等式延拓原理是 \(\hat{\mathbf{K}}\) 上的绝对值。

